% ===================================================================
%  LENTELĖ Nr. 7 — Kaimas žiemą. Sodyba pavasarį.
%  Traduction 2026 d'apres texto/io, controlee sur texto/fr, texto/ru
%  et texto/cs. Consigne : voir texto/lt/10-lentele-01.tex.
%
%  LA FERME EST, POUR LE GALICIEN, L'ENDROIT OU LA DISTANCE A LA
%  LANGUE VOISINE EST MAXIMALE. POUR LE LITUANIEN C'EST DAVANTAGE :
%  c'est l'endroit ou la langue n'a rien a chercher nulle part.
%  Soixante betes, outils et travaux, et pas un emprunt : « karvė »
%  la vache, « jautis » le taureau, « veršis » le veau, « avis » la
%  brebis, « avinas » le belier, « ėriukas » l'agneau, « ožka » la
%  chevre, « kiaulė » la truie, « paršelis » le porcelet, « višta »
%  la poule, « gaidys » le coq, « antis » le canard, « žąsis » l'oie,
%  « triušis » le lapin, « dalgis » la faux, « šakės » la fourche,
%  « grėblys » le rateau, « akėčios » la herse, « šienas » le foin,
%  « vaga » le sillon.
%
%  ET LE TABLEAU 7 RENVERSE CE QUE LES SIX PRECEDENTS AVAIENT
%  ETABLI. Depuis le tableau 6 on sait que le lituanien nomme la
%  chose par l'action : « žadintuvas » ce qui reveille, « keptuvė »
%  ce qui frit, « batsiuvys » celui qui coud les bottes, « kirpėjas »
%  celui qui coupe, « ratininkas » celui de la roue. ICI C'EST
%  L'OUTIL QUI NOMME LA BETE. « arklas » est l'araire, le vieil
%  instrument de bois ; « arklys » le cheval en derive, et veut donc
%  dire l'animal de l'araire. La bete la plus ordinaire de la ferme
%  porte le nom de son travail.
%
%  ON MESURE LA PAR OU LE LITUANIEN A CHANGE. Le meme tableau donne
%  « plūgas », la charrue de fer, qui est un emprunt germanique venu
%  avec l'instrument nouveau. La langue garde donc les deux : le mot
%  ancien qui a fait le nom du cheval, et le mot recent qui n'a fait
%  le nom de rien. Un emprunt n'engendre pas ; un mot du fonds, si.
%
%  LA BLAGUE DU CORDONNIER TIENT EN LITUANIEN, ce qui n'allait pas de
%  soi. Le fac-simile oppose « shuifisto » et « shureparisto »,
%  celui qui fait les chaussures et celui qui les repare. Le
%  lituanien a les deux mots tout faits et batis sur le meme modele —
%  « batsiuvys » celui qui coud les bottes, « batataisys » celui qui
%  les raccommode — de sorte que l'opposition se lit sans glose.
%
%  DEUX GENITIFS TENUS PAR UN AUTRE TOUR, comme aux tableaux 4 et 5 :
%  la litiere (23) de la jument (21) et la niche (38) du chien (37).
%  Une relative garde chaque fois l'ordre de la source.
% ===================================================================

%%K t07-noto-1 noto
\VUnotes{143.2mm}{%
(*) Dėl valgio smulkmenų žiūrėk 6 ir 13 lenteles.%
}

%%K t07-apar-1 apar
\VUsaut{4.0mm}
\VUcentre{13.2pt}{90}{ANTROJI SERIJA}
\VUsaut{3.0mm}
\VUfilet{16mm}
\VUsaut{5.6mm}
\VUtitre{15.0pt}{60}{LENTELĖ Nr. 7}
\VUsaut{3.0mm}

%%K t07-tit-1 sub
\VUcentre{12.0pt}{0}{\textit{Pirmoji scena.}}
\VUsaut{3.0mm}

%%K t07-tit-2 sub
\VUpk{10.2pt}{40}{Kaimas žiemą.}
\VUsaut{4.0mm}

%%K t07-c1-01-1 p
1. --- Per Naujųjų metų atostogas palydėjau tėvą į Naujamiestį, miestelį, kur jam
reikėjo sutvarkyti keletą prekybos reikalų.

%%K t07-c1-02-1 p
2. --- Važiavome \VUgras{pašto karieta} \textsuperscript{(11)}, kuri kursuoja tarp
miesto ir to miestelio; bet kadangi buvo labai šalta, ta kelionė nelabai patogia
priemone nebuvo itin maloni.

%%K t07-c1-03-1 p
3. --- Todėl pajutome tikrą malonumą, kai pamatėme \VUgras{vežėją} sustabdant savo
vežimą aikštėje, prieš \textit{Baltosios Meškos} \VUgras{užeigą} \textsuperscript{(2)}.
\VUgras{Užeigos šeimininkas} \textsuperscript{(4)} laukė mūsų ant savo namo slenksčio,
ramiai rūkydamas savo \VUgras{pypkę.}

%%K t07-c1-03-2 p
Jis pakvietė mus užeiti, ir aš nesileidau ilgai prašomas atsisėsti prieš šakelių
ugnį, kuri traškėjo židinyje. Tada gerokai pasišaipiau iš aštraus \VUgras{šiaurės
vėjo,} kuris girgždina seną \VUgras{iškabą} \textsuperscript{(3)} ir juodais sūkuriais
nunešė \VUgras{dūmus} \textsuperscript{(1)}, kylančius iš židinio.

%%K t07-c1-04-1 p
4. --- Buvo pusryčių metas. Ilgiau nelaukdami, gerai suvalgėme paprastą, bet skanų
patiekalą, kurį mums patiekė (*).

%%K t07-tit-3 sub
\VUpk{10.2pt}{40}{Keli apsilankymai.}
\VUsaut{3.0mm}

%%K t07-c1-05-1 p
5. --- Po pusryčių palydėjau tėvą, kuris yra draudikas, pas jo klientus. Pirmiausia
aplankėme \VUgras{kalvį} \textsuperscript{(5)}, kuris gyvena prie užeigos. Jis buvo
užsiėmęs kaustydamas arklį. Gėrėjausi jo padėjėjo tvirtumu: jis tvirtai laikė ant
savo odinės prijuostės arklio \VUgras{kanopą,} o kalvis stipriais plaktuko smūgiais
tvirtino \VUgras{pasagą} \textsuperscript{(6)}.

%%K t07-c1-06-1 p
6. --- Paskui nuėjome pas kaimo \VUgras{batsiuvį} \textsuperscript{(7)}, senąjį
Jokūbą. Nors jo iškaba yra puikus \VUgras{batas,} jis tenkinosi perpadindamas seną
kiaurų batų porą.

%%K t07-c1-06-2 p
Senis juokdamasis mums pasakė: „Mieste buvau „batsiuvys“ ir neturėjau klientų. Čia
esu „batataisys“, ir darbo daug.“

%%K t07-c1-07-1 p
7. --- Jis papasakojo mums apie savo kaimyną \VUgras{kirpėją} \textsuperscript{(8)},
kurio klientai nepatenkinti. Jo \VUgras{skustuvai} visada išsproginėję, o jo
\VUgras{žirklės} niekada gerai nekerpa.

Bet kadangi jis vienintelis kaimo kirpėjas, savaime suprantama, esi priverstas pas jį
skustis ir kirptis. Be to, jis šykštus ir nemalonus žmogus.

%%K t07-c1-08-1 p
8. --- Laimei, kiti \VUgras{kaimynai} į jį nepanašūs. Antai \VUgras{ratininkas}
\textsuperscript{(9)} yra geras žmogus, darbštus ir paslaugus. Malonu pas jį taisyti
vežimą. Jo darbas visada iki galo atliktas.

%%K t07-c1-09-1 p
9. --- Senasis Jokūbas taip pat nesiskundė \VUgras{balniumi} \textsuperscript{(10)},
kuriam kartais padeda siūti \VUgras{pakinktus,} kai darbas skuba.

%%K t07-c1-09-2 p
Mano tėvas paklausė jį apie šeimą: „Visi sveiki. Mano anūkė \VUgras{mokykloje}
\textsuperscript{(12)}. Jos \VUgras{mokytoja} \textsuperscript{(13)} labai ja
patenkinta, ji gera \VUgras{mokinė} \textsuperscript{(14)}. Ji nepanaši į mano nenaudėlį
sūnų; jis tegalvoja apie žaidimus. \VUgras{Mokytojas} \textsuperscript{(15)}
nepajėgia nieko jo išmokyti.“

%%K t07-tit-4 sub
\VUsaut{3.0mm}
\VUpk{10.2pt}{40}{Kaimo šnekos.}
\VUsaut{3.0mm}

%%K t07-c1-10-1 p
10. --- Paskui senis papasakojo mums kaimo naujienas. Ketinama statyti naują
mokyklą, nes įrengtoji \VUgras{valsčiaus namuose} pasidarė per maža. Beje, tai
lengva, nes pakaks nupirkti dalį \VUgras{malūnininko} sodo. Tai bus labai naudinga
vargšui žmogui. Dėl šalčio \VUgras{malūno} \textsuperscript{(16)} \VUgras{girnos}
nebegali malti kviečių, nes \VUgras{ratą} \textsuperscript{(17)} sustabdė
\VUgras{ledas} \textsuperscript{(25)}, sukaustęs \VUgras{upelį} \textsuperscript{(23)}.

%%K t07-c1-11-1 p
11. --- „Kalbant apie statybas, ar matėte mūsų naują \VUgras{bažnyčią}
\textsuperscript{(18)}? Ji labai vaizdinga po savo sniego skraiste. \VUgras{Varnos}
\textsuperscript{(19)}, kurios nebeatpažįsta savo senosios varpinės, liūdnai tupi ant
didžiojo \VUgras{kapinių} \textsuperscript{(20)} medžio.“

%%K t07-c1-12-1 p
12. --- Ta mintis apie kapines primena batsiuviui žmones, kuriuos mano tėvas buvo
pažinojęs ir kurie mirė pernai. „Kiek \VUgras{kapų} \textsuperscript{(21)}, mano
gerasis pone, prisidėjo prie kitų po \VUgras{kiparisu} \textsuperscript{(22)}! Mirtis
pakirto mūsų senąjį notarą. Visi kaimiečiai dalyvavo jo \VUgras{laidotuvėse.}“

%%K t07-c1-13-1 p
13. --- „Štai jo įpėdinis, kuris čiuožia upeliu. Jis ką tik atėjo, kad pataisyčiau
jo \VUgras{pačiūžos} \textsuperscript{(26)} dirželį, kuris buvo nutrūkęs.“

%%K t07-c1-14-1 p
14. --- „Štai ir ant upelio kranto naujasis \VUgras{laukų sargas}
\textsuperscript{(27)}. Tai buvęs žandaras, kuris gąsdina kaimo padaužas. Jie tuoj
pat pabėga, vos pamatę jo \VUgras{antblauzdžius} \textsuperscript{(29)} ir jo
\VUgras{kepurę} \textsuperscript{(28)}.“

%%K t07-c1-15-1 p
15. --- „A! štai senasis Juozas, kuris veža \VUgras{malkų glėbių vežimėlį}
\textsuperscript{(30)} kepėjui. --- Tiesa, tarė mano tėvas, atpažįstu jo
\VUgras{mulų} \textsuperscript{(31)} pakinkymą. Man kaip tik reikia su juo pasikalbėti,
pasinaudosiu proga. Iki pasimatymo, tėve Jokūbai.“

%%K t07-c1-16-1 p
16. --- Kaip tik kai išėjome, kaimo vaikai paliko mokyklą ir bėgdami puolė prie
\VUgras{čiuožyklos} \textsuperscript{(33)}, rizikuodami parkristi ir susilaužyti
galūnę \VUgras{krisdami} \textsuperscript{(34)}. Vienas iš jų pasiėmė savo
\VUgras{roges} \textsuperscript{(32)} pas ratininką, pasodino jose vieną savo draugą ir
bėgte nudūmė.

%%K t07-c1-17-1 p
17. --- Kiti puolė prie \VUgras{sniego pusnies} \textsuperscript{(37)} netoli
\VUgras{tilto} \textsuperscript{(40)} ir ėmė bombarduoti \VUgras{sniego statulą}
\textsuperscript{(39)}, kurią pastatė vakar dieną ir kuriai užmovė skrybėlę, pypkę ir
mokyklinį krepšį per petį.

%%K t07-c1-18-1 p
18. --- Jiems nelabai rūpi šaldantis vėjas ir šaltis. Daugumas net neturi
\VUgras{šaliko} \textsuperscript{(35)}, ir, kad sušiltų po mūšio \VUgras{sniego
gniūžtėmis} \textsuperscript{(38)}, jiems pakanka saujos labai karštų \VUgras{kaštonų}
\textsuperscript{(42)}, nupirktų iš senosios \VUgras{pardavėjos} Jokūbinos, kuri dreba
nuo šalčio po savo per plona \VUgras{skara} \textsuperscript{(43)}.

%%K t07-tit-5 sub
\VUsaut{3.0mm}
\VUcentre{12.0pt}{0}{\textit{Antroji scena.}}
\VUsaut{3.0mm}

%%K t07-tit-6 sub
\VUpk{10.2pt}{40}{Sodyba pavasarį.}
\VUsaut{4.0mm}

%%K t07-c2-01-1 p
1. --- Nepaisant visų žiemos malonumų, giliai džiaugdamiesi sveikiname
\VUgras{pavasario} sugrįžimą.

%%K t07-c2-01-2 p
Koks linksmas vaizdas yra sodybos kiemas vakare, gegužės mėnesį!

%%K t07-c2-02-1 p
2. --- Ties horizontu \VUgras{saulė} \textsuperscript{(1)} lėtai gula, užliedama
\VUgras{laukus} \textsuperscript{(3)} savo purpuriniais \VUgras{spinduliais}
\textsuperscript{(2)}. Darbštus \VUgras{artojas} \textsuperscript{(6)} skuba suarti
savo \VUgras{dirvą} \textsuperscript{(4)}.

%%K t07-c2-02-2 p
Savo \VUgras{plūgu} \textsuperscript{(5)}, pakinkytu prie stipraus \VUgras{arklio}
\textsuperscript{(7)}, jis brėžia \VUgras{vagą} \textsuperscript{(8)}, į kurią
\VUgras{sėjėjas} \textsuperscript{(9)} pilna sauja mes \VUgras{sėklą,} kuri duos
auksines varpas.

%%K t07-c2-03-1 p
3. --- \VUgras{Šienpjoviai} \textsuperscript{(11)} su savo dalgiais nupjovė pievos
žolę, šiendžiovintojai išdžiovino \VUgras{šieną} \textsuperscript{(10)}, versdami jį
\VUgras{šakėmis} \textsuperscript{(12)} ir grėbliais, ir štai jie grįžta.

%%K t07-c2-03-2 p
Vieni eina prieš sunkų vežimą, traukiamą jaučių; jie neša savo įrankį ant peties.
Kiti, tingesni, patogiai įsitaisė ant kvapnaus šieno.

%%K t07-c2-04-1 p
4. --- Eidami pro šalį jie draugiškai pasveikina \VUgras{sodininką}
\textsuperscript{(16)}, užsiėmusį \VUgras{vaismedžių sode} \textsuperscript{(13)}
genėdamas \VUgras{vaismedžius} ir skiepydamas \VUgras{kriaušes}
\textsuperscript{(14)}. \VUgras{Slyvos} \textsuperscript{(15)} pražydo, o gerai
patręštame \VUgras{darže} \textsuperscript{(17)} daržovės, \VUgras{kopūstai} ir
\VUgras{salotos} jau gerai auga.

%%K t07-c2-04-2 p
Ant mūrelio gražus \VUgras{povas} \textsuperscript{(18)} skleidžia savo uodegą, kad
duotų pasigėrėti jos puikiomis plunksnomis.

%%K t07-tit-7 sub
\VUpk{10.2pt}{40}{Sodybos kiemas.}
\VUsaut{3.0mm}

%%K t07-c2-05-1 p
5. --- \VUgras{Arklidės} \textsuperscript{(19)} durys atvertos, ir matyti ėdžios bei
\VUgras{pakratai} \textsuperscript{(23)}, skirti \VUgras{kumelei}
\textsuperscript{(21)}, kurią \VUgras{arklininkas} \textsuperscript{(20)} ką tik
iškinkė. \VUgras{Kumeliukas} \textsuperscript{(22)}, toks juokingas savo ilgomis
kojomis, šokinėdamas seka paskui motiną.

%%K t07-c2-06-1 p
6. --- Prie \VUgras{šulinio} \textsuperscript{(29)} \VUgras{karvės}
\textsuperscript{(25)} eina gerti į medinį \VUgras{lovį} \textsuperscript{(26)}, kuris
joms tarnauja girdykla. \VUgras{Veršiukas} \textsuperscript{(24)} pasinaudoja proga
pažįsti, o \VUgras{jautis} \textsuperscript{(27)}, uosdamas gerą pašaro kvapą,
džiaugsmingai baubia ir išdidžiai kelia galvą su galingais \VUgras{ragais}
\textsuperscript{(28)}.

%%K t07-c2-07-1 p
7. --- Visi dirba. \VUgras{Tarnaitė} \textsuperscript{(32)} rankoje laiko šulinio
\VUgras{virvę} \textsuperscript{(31)}. Ji semia vandenį \VUgras{kibiru}
\textsuperscript{(30)}, kad pagirdytų gyvulius. Prie \VUgras{kluono}
\textsuperscript{(33)} vyras grėbia šiaudus, o prieš \VUgras{trobos}
\textsuperscript{(34)} duris senoji \VUgras{verpėja} \textsuperscript{(35)} savo
rateliu ir savo \VUgras{verpste} verpia \VUgras{linus} ar \VUgras{kanapes} kaip
senovėje.

%%K t07-tit-8 sub
\VUsaut{3.0mm}
\VUpk{10.2pt}{40}{Ūkininkas.}
\VUsaut{3.0mm}

%%K t07-c2-08-1 p
8. --- \VUgras{Ūkininkas} \textsuperscript{(36)} vadovauja visiems tiems darbams ir
duoda nurodymus savo bernams. Jis liepia pastogėn padėti \VUgras{akėčias}
\textsuperscript{(40)} ir \VUgras{kaplį} \textsuperscript{(39)}, kuriuos pamiršo prie
\VUgras{būdos} \textsuperscript{(38)}, kurioje gyvena \VUgras{šuo}
\textsuperscript{(37)}.

%%K t07-c2-09-1 p
9. --- Jo šūksniai išgąsdina \VUgras{balandį} \textsuperscript{(41)}, kuris visais
sparnais nulekia į \VUgras{balandinę} \textsuperscript{(42)}, ir \VUgras{vištas}
\textsuperscript{(43)}, kurios skuba grįžti į \VUgras{vištidę} \textsuperscript{(44)}.

%%K t07-c2-10-1 p
10. --- Prieš \VUgras{avidę} \textsuperscript{(48)} štai senasis \VUgras{piemuo}
\textsuperscript{(45)}, kuris parveda savo \VUgras{bandą} \textsuperscript{(49)}.
Laikydamas savo \VUgras{lazdą} \textsuperscript{(46)}, kurios jam niekada netrūksta,
kaip ir jo išblukusios \VUgras{palaidinės} \textsuperscript{(47)}, jis veda į tvartą
\VUgras{avinus} \textsuperscript{(50)}, \VUgras{avis} \textsuperscript{(53)},
\VUgras{ėriukus} \textsuperscript{(54)}, \VUgras{ožkas} \textsuperscript{(51)} ir
\VUgras{ožiukus} \textsuperscript{(52)}. Jo ištikimas \VUgras{šuo}
\textsuperscript{(55)} Argas ateina paskutinis ir savo lojimu skatina vėluojančiųjų
skubumą.

%%K t07-c2-11-1 p
11. --- Visas tas triukšmas ir judėjimas netrikdo ramybės geros \VUgras{pieningos
karvės} \textsuperscript{(56)}, kuri skambina savo \VUgras{varpeliu}
\textsuperscript{(57)}, kol ją melžia prieš \VUgras{tvarto} \textsuperscript{(58)}
duris.

%%K t07-c2-12-1 p
12. Ji tarsi supranta, kad turi duoti savo pieno, kad \VUgras{ūkininkė}
\textsuperscript{(60)} galėtų paruošti \VUgras{sviesto} \VUgras{muštuvėje}
\textsuperscript{(61)} arba puikių \VUgras{sūrių} \textsuperscript{(64)}, kurie varva
nuo lentos, padėtos ant \VUgras{kubilo} \textsuperscript{(63)}.

%%K t07-c2-13-1 p
13. --- Visą \VUgras{pieninę} \textsuperscript{(59)} laiko labai švarią \VUgras{ūkio
bernas} \textsuperscript{(65)}, kuris ką tik išplovė \VUgras{pieno puodus}
\textsuperscript{(62)} dideliame kibire, kurį neša rankoje.

%%K t07-tit-9 sub
\VUsaut{3.0mm}
\VUpk{10.2pt}{40}{Paukštidė.}
\VUsaut{3.0mm}

%%K t07-c2-14-1 p
14. --- Su savo storomis medinėmis klumpėmis jis eina kiek nerangiai ir taip
išbaido \VUgras{kalakutą} \textsuperscript{(68)}, \VUgras{antis}
\textsuperscript{(66)}, \VUgras{antinus} \textsuperscript{(67)} ir \VUgras{žąsis}
\textsuperscript{(69)}, kurios plačiai atveria \VUgras{savo snapus}
\textsuperscript{(70)} ir nerimastingai klykia.

%%K t07-c2-14-2 p
\VUgras{Gaidys} \textsuperscript{(71)}, kovingesnis, pastato savo raudoną
\VUgras{skiauterę} \textsuperscript{(72)}, papurto savo \VUgras{uodegos}
\textsuperscript{(73)} plunksnas ir garsiai užgieda „kakariekū“.

%%K t07-c2-15-1 p
15. --- \VUgras{Višta motina} \textsuperscript{(74)} lesinėja
\VUgras{mėšlyne} \textsuperscript{(81)} su savo \VUgras{viščiukais}
\textsuperscript{(75)}. \VUgras{Paršelis} \textsuperscript{(78)} bėga pas savo motiną,
milžinišką \VUgras{kiaulę} \textsuperscript{(77)}, kurią matome prie kiaulidės.

%%K t07-c2-16-1 p
16. --- Savo gardeliuose \VUgras{triušiai} \textsuperscript{(79)} godžiai ėda
gležną \VUgras{žolę} \textsuperscript{(80)}, kurią jiems atneša ūkininko duktė.
Jonelė labai myli savo triušius ir kasdien, pasiėmusi šviežius \VUgras{kiaušinius}
\textsuperscript{(76)} vištidėje, ateina aplankyti savo draugužių.

\VUsaut{6.0mm}
\VUfilet{22mm}
